\documentclass[11pt]{article}
\usepackage[margin=1in]{geometry}
\usepackage[T1]{fontenc}
\usepackage[utf8]{inputenc}
\usepackage{amsmath,amssymb,amsthm}
\usepackage{graphicx}
\usepackage{xcolor}
\usepackage{microtype}
\usepackage[colorlinks=true,linkcolor=blue!55!black,citecolor=blue!55!black,urlcolor=blue!55!black]{hyperref}

\newtheorem{theorem}{Theorem}
\newtheorem{corollary}{Corollary}
\newcommand{\Zp}{\mathbb Z_+}
\newcommand{\Ap}{\mathcal A}

\title{The discrete maximal operator blocks direct extrapolation below one}
\author{Candidate full-counterexample packet for arXiv:2304.14250}
\date{11 August 2026}

\begin{document}
\maketitle

\begin{center}
\fbox{\parbox{0.92\textwidth}{\textbf{Status: candidate full counterexample to the direct interpretation, likely valid, subject to human review.}
The discrete Hardy--Littlewood maximal operator satisfies the source's
weighted hypothesis at every fixed $p_0>1$, but is not bounded on unweighted
$\ell^p(\Zp)$ for any $0<p\le1$.  Therefore no below-one extension of the
source theorem that includes the constant weight can be true under the same
hypothesis.}}
\end{center}

\section{The signal and its necessary interpretation}

Saker, Saied, and Agarwal prove discrete Rubio de Francia extrapolation for
$1<p<\infty$ and say in the abstract and Introduction that the case
$0<p<1$ is open \cite{source}.  Their discrete Muckenhoupt class $\Ap_p$ is
defined only for $p>1$, so a literal conclusion with ``$w\in\Ap_p$'' is not
defined below one.  Nevertheless, every direct extension of their operator
theorem has one unavoidable consequence: it must give the unweighted
$\ell^p$ estimate, since the constant weight is the baseline weight at every
classical exponent.

That necessary consequence is false.  Thus the result below answers the
only direct interpretation of the printed signal.  It also diagnoses why a
valid below-one theorem must change either the initial hypothesis, the
target weight class, the sequence family, or the target norm.

\section*{Idea of the proof}

Use the operator that generates the weight classes in the first place.  The
discrete Hardy--Littlewood maximal operator obeys the required weighted
strong-type estimate for every $p_0>1$ and every $\Ap_{p_0}$ weight.  But on
the point mass $e_1$, the interval $\{1,\ldots,n\}$ gives
$Me_1(n)\ge1/n$.  This tail is outside $\ell^p$ exactly when $p\le1$.

\section{The counterexample}

For a sequence $f$ on $\Zp=\{1,2,\ldots\}$, let
\[
 Mf(n)=\sup_{\substack{I\subset\Zp\text{ finite interval}\\ n\in I}}
 \frac1{|I|}\sum_{k\in I}|f(k)|
\]
be the discrete uncentered Hardy--Littlewood maximal operator.  The same
argument works for any one-sided variant which takes the initial intervals
$\{1,\ldots,n\}$ among its averages.

\begin{theorem}\label{thm:counterexample}
Fix $p_0>1$.  The operator $M$ satisfies the source hypothesis
\begin{equation}\label{eq:weighted}
 \sum_{n=1}^{\infty}(Mf(n))^{p_0}w(n)
 \le C\bigl(p_0,[w]_{\Ap_{p_0}}\bigr)
 \sum_{n=1}^{\infty}|f(n)|^{p_0}w(n)
\end{equation}
for every discrete $\Ap_{p_0}$ weight $w$.  However, for every
$0<p\le1$, there is no constant $C_p$ such that
\[
 \sum_{n=1}^{\infty}(Mf(n))^p
 \le C_p\sum_{n=1}^{\infty}|f(n)|^p
\]
for all finitely supported $f$.
\end{theorem}

\begin{proof}
The weighted strong-type estimate \eqref{eq:weighted} is the defining
maximal theorem for discrete Muckenhoupt weights.  It is recalled explicitly
in the source and proved in the discrete setting by Saker and Mahmoud
\cite{sakermahmoud}.  Replacing its dependence on the characteristic by a
monotone envelope, if necessary, gives exactly the positive increasing
control function required by the source extrapolation hypothesis.

Let $e_1$ be the point mass at $1$.  For every $n\ge1$, the interval
$I_n=\{1,\ldots,n\}$ contains $n$, and hence
\[
 Me_1(n)\ge \frac1n\sum_{k=1}^n e_1(k)=\frac1n.
\]
Consequently, for $0<p\le1$,
\[
 \sum_{n=1}^{\infty}(Me_1(n))^p
 \ge \sum_{n=1}^{\infty}\frac1{n^p}=\infty,
 \qquad
 \sum_{n=1}^{\infty}|e_1(n)|^p=1.
\]
Thus $M$ is not bounded on unweighted $\ell^p(\Zp)$ at or below one,
although it satisfies the full weighted $p_0$ hypothesis.
\end{proof}

\begin{corollary}[family-of-pairs form]
Let
\[
 \Omega=\{(Mf,f): f\ge0\text{ and finitely supported on }\Zp\}.
\]
At exponent $p_0>1$, the source's weighted inequality holds for every pair
in $\Omega$ and every $w\in\Ap_{p_0}$.  For every $0<p\le1$, its
unweighted conclusion fails for the pair $(Me_1,e_1)$.
\end{corollary}

This removes a possible operator-linearity loophole: the source emphasizes
that its theorem is about arbitrary families of pairs and does not require
linearity or sublinearity.  The counterexample already works in that
generality.

\section{Scope, stronger variants, and novelty}

The counterexample does not contradict the source's final $\Ap_\infty$
corollary.  That corollary assumes the initial $p_0$ inequality for
\emph{every} $\Ap_\infty$ weight, a strictly stronger premise than the main
$\Ap_{p_0}$ hypothesis.  Nor does it contradict extrapolation for decreasing
or quasi-nonincreasing sequences with Ari\~no--Muckenhoupt-type weight
classes.  A 2026 paper proves an upward extrapolation theorem for a
restricted quasi-nonincreasing family and a different
$\mathcal{QB}_{\beta,p}$ class \cite{singh}; it does not supply the direct
all-sequence $\Ap_p$ extension refuted here.

Accordingly, a meaningful below-one replacement must be reformulated.  The
source's own $\Ap_\infty$ corollary is one option; restricted monotonicity,
$B_p$-type classes, weak-type targets, or Hardy-space substitutes are other
possible directions.  The present theorem rules out retaining only the main
$\Ap_{p_0}$ premise while demanding an $\ell^p$ conclusion that includes
the constant weight.

On 11 August 2026, the run registry and all cheap indexes had no packet for
arXiv:2304.14250.  Exact local-corpus and bounded primary-source searches
found the source and the later restricted-sequence paper, but no resolution
of this exact signal.  This is a bounded novelty check, not an exhaustive
priority claim.  The obstruction itself is elementary and should be treated
as a clarification/counterexample rather than a claim of deep novelty.

\appendix
\section{Source-page evidence}

Figure~\ref{fig:source} reproduces the source abstract, including the
below-one open signal.

\begin{figure}[p]
\centering
\includegraphics[height=0.89\textheight]{figures/open_problem_page.png}
\caption{The first printed source page: the abstract states that the case
$0<p<1$ is an open problem.}
\label{fig:source}
\end{figure}

\begin{thebibliography}{9}
\bibitem{source}
S. H. Saker, A. I. Saied, and R. P. Agarwal,
\emph{Discrete Rubio de Francia extrapolation theorem via factorization of
weights and iterated algorithms}, arXiv:2304.14250 [math.FA], 2023.

\bibitem{sakermahmoud}
S. H. Saker and R. R. Mahmoud,
\emph{Boundedness of both discrete Hardy and Hardy--Littlewood maximal
operators via Muckenhoupt weights}, Rocky Mountain J. Math. 51 (2021),
733--746.

\bibitem{singh}
M. Singh, A. Gogatishvili, R. Panchal, and A. P. Singh,
\emph{Rubio de Francia extrapolation theorem for quasi non-increasing
sequences}, arXiv:2603.06109 [math.FA], 2026.
\end{thebibliography}

\end{document}
